% Generated draft: requires mathematical and source review.
\documentclass[11pt]{article}
\usepackage[margin=25mm]{geometry}
\usepackage{fontspec}
\setmainfont{Latin Modern Roman}
\usepackage{amsmath,amssymb,mathtools,unicode-math}
\setmathfont{Latin Modern Math}
\usepackage{enumitem,needspace,xurl,hyperref}
\setlength{\parindent}{0pt}
\setlength{\parskip}{5pt}
\setlength{\emergencystretch}{3em}
\newcommand{\N}{\mathbb{N}}
\newcommand{\Z}{\mathbb{Z}}
\newcommand{\Q}{\mathbb{Q}}
\newcommand{\R}{\mathbb{R}}
\newcommand{\C}{\mathbb{C}}
\newcommand{\F}{\mathbb{F}}
\newcommand{\A}{\mathbb{A}}
\newcommand{\PP}{\mathbb{P}}
\newcommand{\E}{\mathbb{E}}
\newcommand{\eps}{\varepsilon}
\newcommand{\dd}{\,\mathrm{d}}
\begin{document}
\section*{UnsolvedMath ID: 20002379}
\textbf{Source problem code:} catalogue entry
\subsection{Definitions and mathematical statement}
\textbf{Participant context.} A.7 is Gunter Fuchs's 2004 contribution describing interests in the method refuting the Cofinal Branches Hypothesis (CBH), and in inner models built from extender sequences together with partial iteration strategies and broadly applicable fine structure. It states no formal open problem, specified tree or universal fine-structure theorem. The following is an auxiliary conditional synthesis from the aid.

Fix an iteration tree $\mathcal T$ of limit length $\delta$ in a specified iteration-theoretic setting. Assume it has no cofinal branch with an externally well-founded direct-limit model, while at every proper limit stage $\beta<\delta$ its actual continuation branch $b_\beta$ through $\mathcal T\restriction\beta$ has a well-founded branch model. A branch-correct partial choice along this tree assigns $b_\beta$ to an admitted proper limit restriction. This is not automatically a full winning iteration strategy on all possible trees.

\textbf{Supplied final-gap obstruction.} Such choices can be made on all proper limit restrictions, but no branch-correct extension can add $\mathcal T$ itself to the domain. If the domain is closed under taking initial limit restrictions, its admitted indices are downward closed. Thus omission of one proper limit restriction forces omission of every later proper limit restriction. It either admits them all with only a final gap, or has a tail gap among those indices.

To conclude that directed-union closure rules out the single final gap, the proper limit restrictions must form a directed cofinal family with union $\mathcal T$. Under that additional cofinality hypothesis, admitting all of them and closing under their union would force admission of $\mathcal T$, a contradiction. Therefore a branch-correct initial-restriction-closed, directed-union-closed domain omits a cofinal tail of proper limit restrictions. This cofinality condition matters: a general limit ordinal need not be a limit of proper limit stages.

\textbf{Supplied copying obstruction.} Suppose a copied tree $\mathcal U$ comes with a bijection between its cofinal branches and those of $\mathcal T$, and for each corresponding pair a membership-preserving embedding from the original branch model into the copied one. If an original branch model has an external descending membership sequence, its image is a descending sequence in the copied model. Thus every copied cofinal branch model is ill founded as well; the copied tree remains a CBH counterexample. Branch bijection alone, without the model maps preserving the relevant membership relation, would not establish this conclusion.

The aid cites a Neeman--Steel example of length $\omega\cdot\omega=\omega^2$. Its proper positive limit stages are exactly $\omega n$, $1\le n<\omega$, so they are cofinal in its length. It reports that Corollary 3.8 and a Martin--Steel input verify the well-founded actual continuations, while a final critical-point sequence gives the ill-founded union. The paper, complete assumptions and critical-point construction are not supplied here. These are attributed conditional reports, not a new construction or an unconditional theorem about every iteration tree. They illustrate the participant's two interests without turning his note into a solved problem.
\subsection{Short English statement}
Fuchs's entry is a statement of workshop interests. The aid's conditional example separates locally well-founded continuation choices from a final CBH obstruction. Initial-restriction and directed-union closure force a tail gap when proper limit stages are cofinal, and faithful branch-model copying preserves the ill-foundedness obstruction.
\subsection{Sources}
\begin{itemize}
\item Local UnsolvedMath registry, permanent ID 20002379, source code catalogue entry; canonical \texttt{aim-logic-notes.json}, notes[154], statement status reconstructed unverified.
\item AIM 2004 \emph{Recent advances in core model theory}, participant contribution A.7, Gunter Fuchs: \url{https://aimath.org/WWN/coremodel/coremodel.pdf}. The record reports official-PDF confirmation, not performed here.
\item The aid names Neeman--Steel, Corollary 3.8, and a Martin--Steel input, without full paper details, theorem hypotheses or the final critical-point argument.
\item Registry research attempt 1, dated 2026-08-14: context only, status at run not a problem, rigor ``proved'', novelty confidence low. The explicitly machine-generated, non-peer-reviewed aid requires independent verification before citation.
\end{itemize}
\end{document}
